% Appendix: Repository Structure
\chapter*{Appendix A: Project Repository Structure}
\addcontentsline{toc}{chapter}{Appendix A: Project Repository Structure}

This appendix provides a detailed overview of the project repository organization, including all directories, notebooks, scripts, and datasets used throughout the analysis.

\section*{Repository Overview}
The project repository is organized into five main directories, each serving a specific purpose in the data mining pipeline:

\begin{itemize}
    \item \texttt{analysis/} - Jupyter notebooks for exploratory data analysis and modeling
    \item \texttt{dataset/} - Raw and processed datasets (excluded from version control)
    \item \texttt{models/} - Custom implementations of machine learning algorithms
    \item \texttt{results/} - Generated outputs, processed datasets, and evaluation reports
    \item \texttt{scripts/} - Utility scripts for data download and verification
\end{itemize}

\section*{Complete Directory Structure}

\begin{verbatim}
project-root/
|-- analysis/                    # Jupyter notebooks for data analysis
|   |-- climate.ipynb           # Climate data processing & EDA
|   |-- dbscan.ipynb            # DBSCAN clustering analysis
|   |-- decision_tree.ipynb     # Decision Tree model analysis
|   |-- elevation.ipynb         # Elevation feature extraction
|   |-- feature_engineering.ipynb # Feature engineering pipeline
|   |-- fire.ipynb              # Fire dataset processing
|   |-- kmeans.ipynb            # K-Means clustering analysis
|   |-- knn.ipynb               # KNN model analysis
|   |-- landcover.ipynb         # Land cover analysis
|   |-- merge.ipynb             # Dataset integration
|   |-- model_training.ipynb    # Model training & evaluation
|   |-- random_forest.ipynb     # Random Forest model analysis
|   |-- soil.ipynb              # Soil properties extraction
|   `-- undersampling.ipynb     # Data balancing techniques
|
|-- dataset/                     # Raw datasets (gitignored - large files)
|   |-- climate_dataset/        # CHIRPS & TerraClimate rasters
|   |   |-- prec/               # Precipitation data
|   |   |-- tmax/               # Max temperature data
|   |   `-- tmin/               # Min temperature data
|   |-- elevation_dataset/      # GMTED2010 elevation data
|   |   |-- be15_grd/           # Raw elevation grid
|   |   `-- elevation_grd/      # Processed elevation grid
|   |-- fire_dataset/           # FIRMS fire points (CSV)
|   |   |-- algeria.csv
|   |   |-- fires.csv
|   |   `-- tunisia.csv
|   |-- land_cover_dataset/     # FAO land cover shapefiles
|   |   |-- algeria/
|   |   `-- tunisia/
|   `-- soil_dataset/           # HWSD v2.0 soil data
|       |-- HWSD2.bil
|       |-- HWSD2.hdr
|       |-- HWSD2.mdb
|       |-- HWSD2.prj
|       `-- HWSD2.stx
|
|-- models/                      # Custom ML implementations
|   |-- __init__.py
|   |-- dbscan.py               # DBSCAN clustering algorithm
|   |-- decision_tree.py        # Decision Tree (CART algorithm)
|   |-- kmeans.py               # K-Means clustering algorithm
|   |-- knn.py                  # K-Nearest Neighbors classifier
|   `-- random_forest.py        # Random Forest ensemble
|
|-- results/                     # Output datasets & reports
|
|-- scripts/                     # Utility scripts
|   |-- download_data.py        # Dataset download helper
|   `-- verify_data.py          # Dataset verification tool
|
|-- .gitignore                   # Git ignore rules
`-- requirements.txt             # Python dependencies
\end{verbatim}

\section*{Directory Descriptions}

\subsection*{Analysis Directory}
The \texttt{analysis/} directory contains all Jupyter notebooks used for data exploration, preprocessing, and modeling. Each notebook corresponds to a specific phase of the project:

\begin{itemize}
    \item \textbf{Data Processing Notebooks}: \texttt{climate.ipynb}, \texttt{elevation.ipynb}, \texttt{fire.ipynb}, \texttt{landcover.ipynb}, \texttt{soil.ipynb} handle individual dataset preprocessing and exploratory analysis
    \item \textbf{Integration Notebooks}: \texttt{merge.ipynb} combines all processed datasets, \texttt{feature\_engineering.ipynb} creates derived features, \texttt{undersampling.ipynb} addresses class imbalance
    \item \textbf{Supervised Learning Notebooks}: \texttt{knn.ipynb}, \texttt{decision\_tree.ipynb}, \texttt{random\_forest.ipynb} implement and evaluate classification models
    \item \textbf{Unsupervised Learning Notebooks}: \texttt{kmeans.ipynb}, \texttt{dbscan.ipynb} perform clustering analysis
    \item \textbf{Model Training}: \texttt{model\_training.ipynb} provides comprehensive model comparison and evaluation
\end{itemize}

\subsection*{Dataset Directory}
The \texttt{dataset/} directory stores all raw and intermediate datasets. Due to large file sizes, this directory is excluded from version control via \texttt{.gitignore}:

\begin{itemize}
    \item \textbf{climate\_dataset/}: Contains CHIRPS precipitation data and TerraClimate temperature rasters organized by variable type
    \item \textbf{elevation\_dataset/}: Stores GMTED2010 digital elevation model data in both raw and processed formats
    \item \textbf{fire\_dataset/}: Contains FIRMS (Fire Information for Resource Management System) fire detection data in CSV format for Algeria and Tunisia
    \item \textbf{land\_cover\_dataset/}: Includes FAO land cover classification shapefiles for both countries
    \item \textbf{soil\_dataset/}: Contains Harmonized World Soil Database (HWSD) v2.0 raster and database files
\end{itemize}

\subsection*{Models Directory}
The \texttt{models/} directory contains custom Python implementations of machine learning algorithms:

\begin{itemize}
    \item \textbf{Supervised Learning}: \texttt{knn.py}, \texttt{decision\_tree.py}, \texttt{random\_forest.py} implement classification algorithms from scratch
    \item \textbf{Unsupervised Learning}: \texttt{kmeans.py}, \texttt{dbscan.py} implement clustering algorithms
    \item All implementations follow scikit-learn API conventions for consistency
\end{itemize}

\subsection*{Results Directory}
The \texttt{results/} directory stores all outputs generated during the analysis, including processed datasets, model predictions, evaluation metrics, and visualization outputs.

\subsection*{Scripts Directory}
The \texttt{scripts/} directory contains utility scripts for data management:

\begin{itemize}
    \item \texttt{download\_data.py}: Automates the download of publicly available datasets
    \item \texttt{verify\_data.py}: Validates dataset integrity and completeness
\end{itemize}

\section*{Dependencies}
All required Python packages are specified in \texttt{requirements.txt}, including libraries for data manipulation (pandas, numpy), geospatial analysis (geopandas, rasterio), machine learning (scikit-learn), and visualization (matplotlib, seaborn).